\SourceLeaf{128}
\begin{center}
\textbf{REZENSION.}
\end{center}

\medskip
\noindent\rule{\linewidth}{.3pt}

\noindent Allgemeine Literatur-Zeitung von 1808. Halle-Leipzig 1808. Nr.\,45. 12. Februar. S.\,353--358.

\noindent\rule{\linewidth}{.3pt}
\medskip

Dresden, in der Waltherschen Hofbuchhandlung: \textsc{Leonellis} \emph{logarithmische Supplemente, als ein Beitrag, Mängel der gewöhnlichen Logarithmentafeln zu ersetzen. Aus dem Französischen mit einigen Zusätzen von} \textsc{Gottfried Wilhelm Leonhardi}, \emph{Unterleutnant im kurfürstlich-sächsischen Feldartilleriekorps.} 1806. 88 S. im Oktavformat.

Diese kleine Schrift enthält zwei voneinander unabhängige Abhandlungen mit einem gemeinsamen Zweck. Beide sollen etwas leisten, was mit den gewöhnlichen Logarithmentafeln ohne weitere Hilfsmittel entweder weniger genau oder weniger bequem zu erreichen ist. Die erste soll dazu dienen, mithilfe einiger leicht einsetzbarer, hier mitgelieferter Hilfstafeln den Logarithmus jeder gegebenen Zahl ohne allzu große Mühe auf mehr Dezimalstellen zu berechnen, als die gewöhnlichen Tafeln erlauben. Die andere entwickelt die Idee einer besonderen Tafel, mit der die Logarithmen von Summen oder Differenzen zweier nur durch ihre Logarithmen gegebener Größen in einem einzigen Rechenschritt bestimmt werden können. Das soll bequemer sein als bei anderen Verfahren und nach dem Plan des Verfassers ebenfalls eine ungewöhnlich große Zahl von Dezimalstellen (14) liefern. Von dieser Tafel ist jedoch bisher nur eine Probe beigefügt. Da die numerische Arbeit mit Logarithmen und jede Erleichterung dabei für alle, die viel mit Zahlen rechnen, von großer Bedeutung sind, dürfte es sich lohnen, diesen Untersuchungen eine nicht
